\chapter{Symbols and register ordering}

\section{Symbol table}
The symbols used throughout the thesis are collected below for quick reference:
\begin{table}[htbp]
\centering
\caption{Frequently used symbols.}
\begin{tabular}{ll}
\toprule
Symbol & Meaning \\
\midrule
$N$ & Linear image side length ($4$, $8$, or $16$) \\
$m$ & $2\log_2 N$ position qubits \\
$N_{\mathrm{pix}}$ & $N^2=2^m$ pixels \\
$\theta_p$ & FRQI angle for pixel $p$ \\
$p_1$ & Probability of outcome $\ket{1}$ on the color qubit \\
$A$ & Readout confusion matrix ($2\times 2$) \\
$\Delta$ & Paired difference (vchain$-$naive or mitigated$-$raw) \\
\bottomrule
\end{tabular}
\end{table}

\section{Register ordering}
The implementation fixes the color qubit as the least significant bit of the joint index $(\mathrm{position}\Vert c)$, matching Qiskit bit-ordering in file \texttt{src/frqi.py} and the structural code in file \texttt{src/improved.py}.

\chapter{Reproduction checklist}

Run from the repository root, with the virtual environment set from file \texttt{requirements.txt}:
\begin{enumerate}
  \item \texttt{pytest}
  \item \texttt{python scripts/frqi\_qhed\_sobel.py}
  \item \texttt{python scripts/frqi\_structural\_resources.py}
  \item \texttt{python scripts/noisy\_frqi\_sweep.py} (add \texttt{---allow-heavy-dm} flag only on machines with sufficient RAM)
  \item \texttt{python scripts/noisy\_recon\_qhed\_edge\_metrics.py}
  \item \texttt{python scripts/readout\_mitigation\_shot\_sweep.py}
  \item \texttt{python scripts/build\_results\_tables.py ---require-figures}
\end{enumerate}
See file \texttt{thesis/results\_configuration.md} for manifest timestamps and pinned versions.

\chapter{Supplementary tables}

\Cref{tab:noisy-wide-excerpt} presents a couple of noisy edge metrics, with one row per image size.
\begin{table}[htbp]
\centering
\caption{Noisy edge metrics.}
\label{tab:noisy-wide-excerpt}
\footnotesize
\begin{tabular}{lrrrr}
\toprule
image & scale & edge\_SSIM naive & edge\_SSIM vchain & $\Delta$ \\
\midrule
test\_4x4 & 0.05 & 0.195 & -0.141 & -0.336 \\
test\_8x8 & 0.20 & $\approx 0$ & 0.020 & 0.020 \\
test\_16x16 & 0.10 & 0.037 & --- & --- \\
\bottomrule
\end{tabular}
\end{table}
The complete paired naive-versus-v-chain pivot can be found in file \texttt{thesis/generated/tables.md}; This table reproduces only a small slice. The full CSV can be consulted in \texttt{outputs/noisy\_recon\_qhed\_edges.csv}. Rows where v-chain is missing for image of size $16\times16$ appear as \texttt{NaN} in the CSV when density-matrix simulation is skipped due to memory constraints.

\chapter{Complexity and scaling reference}

Table bellow fixes the orders of magnitude and scalings that appear in this thesis.
\begin{table}[htbp]
\centering
\caption{Orders of magnitude and scaling used in this thesis.}
\begin{tabular}{ll}
\toprule
Object & Scaling \\
\midrule
Statevector amplitudes & $O(2^n)$ complex numbers \\
Density matrix entries & $O(4^n)$ real numbers \\
FRQI qubits & $n = 1 + m$ with $m=2\log_2 N$ \\
Pixels & $N_{\mathrm{pix}} = 2^m$ \\
\bottomrule
\end{tabular}
\end{table}

\chapter{Extended noisy edge metrics}

\Cref{tab:noisy-long} lists additional rows from file \texttt{outputs/noisy\_recon\_qhed\_edges.csv} for \texttt{test\_4x4} and \texttt{test\_8x8} at all recorded noise scales (rounded values).
\begin{table}[htbp]
\centering
\footnotesize
\caption{Noisy reconstruction edge SSIM naive vs v-chain for smaller test images.}
\label{tab:noisy-long}
\begin{tabular}{p{2.0cm}rrrr}
\toprule
image & scale & edge\_SSIM naive & edge\_SSIM vchain & $\Delta$\,edge\_SSIM \\
\midrule
test\_4x4 & 0.00 & 0.196 & $-0.108$ & $-0.303$ \\
test\_4x4 & 0.05 & 0.195 & $-0.141$ & $-0.336$ \\
test\_4x4 & 0.10 & 0.167 & $-0.141$ & $-0.308$ \\
test\_4x4 & 0.15 & 0.195 & 0.219 & 0.024 \\
test\_4x4 & 0.20 & 0.580 & 0.169 & $-0.410$ \\
test\_8x8 & 0.00 & 0.163 & 0.032 & $-0.131$ \\
test\_8x8 & 0.05 & $-0.005$ & $-0.006$ & $-0.000$ \\
test\_8x8 & 0.10 & $0.000$ & $-0.095$ & $-0.095$ \\
test\_8x8 & 0.15 & $0.000$ & $-0.272$ & $-0.272$ \\
test\_8x8 & 0.20 & $0.000$ & 0.020 & 0.020 \\
\bottomrule
\end{tabular}
\end{table}

\chapter{Acronyms and repository naming}

\begin{description}[style=nextline, leftmargin=2.2cm, labelwidth=2cm, font=\normalfont\bfseries]
  \item[FRQI] Flexible representation of quantum images; intensity encoded as $R_y$ angles on a color qubit indexed by a position register.
  \item[QHED] Quantum Hadamard edge detection; Walsh--Hadamard style processing followed by classical Sobel scoring in this repository.
  \item[NISQ] Noisy intermediate-scale quantum; devices large enough to explore heuristics yet dominated by hardware imperfections.
  \item[CX] Controlled-NOT (CNOT) count after transpilation to the native two-qubit gate.
  \item[DM] Density-matrix simulation method in Qiskit Aer (\texttt{density\_matrix}).
  \item[SPAM] State preparation and measurement errors collectively.
  \item[SSIM / PSNR] Structural similarity index and peak signal-to-noise ratio, computed with scikit-image conventions on \texttt{uint8} rasters.
  \item[\texttt{noancilla}] Qiskit \texttt{MCXGate} synthesis mode without dedicated ancillas for multi-controlled $X$.
  \item[\texttt{v-chain}] Qiskit \texttt{MCXGate} dirty-ancilla v-chain decomposition lowering CX depth in many benchmarks.
  \item[\texttt{struct\_naive\_slice\_scaled}] Engineering extrapolation that transpiles one representative naive pixel slice and multiplies depth/CX by $N_{\mathrm{pix}}$.
\end{description}

\chapter{Structural metrics}

\Cref{tab:struct-csv} was created from file \texttt{outputs/frqi\_structural\_metrics.csv}. The values are gathered from \cref{ch:methods}; rows labelled \texttt{struct\_naive\_full} exist only for the image of size $4\times 4$, because full naive transpilation becomes difficult for a larger $m$.
\begin{table}[htbp]
\centering
\caption{Structural FRQI preparation metrics.}
\label{tab:struct-csv}
\scriptsize
\begin{tabular}{@{}lrrlrrr@{}}
\toprule
image & $N$ & $m$ & kind & qubits & depth & CX \\
\midrule
test\_4x4 & 4 & 4 & struct\_naive\_full & 6 & 2436 & 1160 \\
test\_4x4 & 4 & 4 & struct\_naive\_slice & 6 & 162 & 74 \\
test\_4x4 & 4 & 4 & struct\_naive\_slice\_scaled & 6 & 2592 & 1184 \\
test\_4x4 & 4 & 4 & struct\_vchain & 8 & 839 & 331 \\
test\_8x8 & 8 & 6 & struct\_naive\_slice & 8 & 421 & 222 \\
test\_8x8 & 8 & 6 & struct\_naive\_slice\_scaled & 8 & 26944 & 14208 \\
test\_8x8 & 8 & 6 & struct\_vchain & 12 & 5423 & 2147 \\
test\_16x16 & 16 & 8 & struct\_naive\_slice & 10 & 897 & 446 \\
test\_16x16 & 16 & 8 & struct\_naive\_slice\_scaled & 10 & 229632 & 114176 \\
test\_16x16 & 16 & 8 & struct\_vchain & 16 & 29435 & 11663 \\
\bottomrule
\end{tabular}
\end{table}

\chapter{Time running the simulations for $4\times4$ and $8\times8$ images}

\Cref{tab:wallclock} records how long some simulations took to complete.
\begin{table}[htbp]
\centering
\caption{Time needed for running simulations for selected noisy FRQI experiments.}
\label{tab:wallclock}
\footnotesize
\begin{tabular}{@{}llrrr@{}}
\toprule
image & method & scale & time (s) & total qubits \\
\midrule
test\_4x4 & naive & 0.00 & 0.67 & 6 \\
test\_4x4 & vchain & 0.00 & 0.36 & 8 \\
test\_4x4 & naive & 0.20 & 0.82 & 6 \\
test\_4x4 & vchain & 0.20 & 0.34 & 8 \\
test\_8x8 & naive & 0.00 & 11.92 & 8 \\
test\_8x8 & vchain & 0.00 & 49.22 & 12 \\
test\_8x8 & naive & 0.20 & 12.48 & 8 \\
test\_8x8 & vchain & 0.20 & 54.81 & 12 \\
\bottomrule
\end{tabular}
\end{table}

In the case of the $8\times8$ image, v-chain circuits are wider; even with fewer CX than naive circuits after transpilation, scheduling a larger Hilbert space can raise the time for generatiing the density-matrix.

\chapter{Repository map and software architecture}
\label{sec:repo-map}

The complete source code is available at \url{https://github.com/AlexandruAlex10/qip-frqi-qhed-mcx-ancilla-and-calibration-matrix-inversion}.

\Cref{tab:repo-map} ties the mathematical objects of the methods chapter to Python modules and scripts.

\begin{table}[H]
\centering
\caption{Core modules and experiment drivers.}
\label{tab:repo-map}
\small
\begin{tabular}{@{}p{4.2cm}p{4.8cm}p{4.6cm}@{}}
\toprule
\textbf{Concern} & \textbf{Module(s)} & \textbf{Script(s)} \\
\midrule
FRQI angles, address loop & \path{src/frqi.py} & \path{frqi_qhed_sobel.py} \\
Structural \texttt{MCXGate} modes & \path{src/improved.py} & \path{frqi_structural_resources.py} \\
QHED + Sobel scoring & \path{src/qhed.py} & \path{frqi_qhed_sobel.py}, \path{noisy_recon_qhed_edge_metrics.py} \\
Metrics (PSNR, SSIM, fidelity) & \path{src/metrics.py} & all image-quality paths \\
Depolarizing DM sweeps & \path{src/noise_models.py} & \path{noisy_frqi_sweep.py} \\
Readout matrix inversion & (design in \path{thesis/design/}) & \path{readout_mitigation_shot_sweep.py} \\
\bottomrule
\end{tabular}
\end{table}

\section{Computational workflow and software architecture}
The experiments are orchestrated as Python scripts that write CSV outputs and PNG figures under \texttt{outputs/}. The thesis ingestion script (\texttt{build\_results\_tables.py}) validates column contracts, compares CSV modification times against manifest timestamps, and emits Markdown and \LaTeX{} fragments under \texttt{thesis/generated/}. Committed excerpts under \texttt{thesis/masters\_en/generated/} keep the PDF buildable even when the gitignored generated files are absent on a fresh clone.
